\scanpage{186}

Ainsi $P(x)=1$, ce qui implique que $\unsure{Z|\mathfrak z}=k$, donc

\uline{Prop} \deleted{\unsure{toute|Toute}}\inserted{Toute} représentation \inserted{lin.} rationnelle irréductible de $k$ est de dim 1.

\uline{Corollaire} \unsure{Une|Toute} représentation lin. \inserted{rationnelle} de dim finie de $k$ peut se mettre sous forme \deleted{diagonale}\inserted{triangulaire}, avec des $1$ sur la diagonale.

Si on est en caractéristique $p$\inserted{$\ne 0$}, on a \deleted{\unsure{une}} \deleted{\unsure{automorphisme}}\inserted{endomorphisme} \unsure{rationnel|rationnelle|rationnelles} de $k$ \deleted{du type} $F$, $x\unsure{\to|-|\mapsto}x^p$, \deleted{\unsure{qui}} d'où une opération naturelle sur les classes de représentations linéaires, en composant avec $F$ et ses itérés ; ce phénomène est tout à fait distinct de ce qui se passe en car. 0 ; notons que parmi les repr. de dim 2 en car. $p$, on a \unsure{ceux} \deleted{\unsure{qui}} par \deleted{\unsure{rotations type}} translation dans l'espace des polynômes $f(\lambda,\mu)\unsure{:|=|(x)=}\lambda X^{p^h}+\mu$ ($h$ fixé, $X$ l'indéterminée) :

\[\begin{aligned}\tau_y f(\lambda,\mu)\unsure{(X)|\in X)|(x)} &= \lambda(X+y)^{p^h}+\mu\\ &=\lambda X^{p^h}+\mu+\lambda y^{p^h}\unsure{,=|=}f(\lambda,\mu+\lambda y^{p^h}),\quad \text{i.e.}\end{aligned}\]

la représentation $E_h(y)=\begin{pmatrix}1&y^{p^h}\\0&1\end{pmatrix}$ ;

les représentations, pour \deleted{\unsure{des}} \deleted{\unsure{variables}} $h$ distincts, sont inéquivalentes deux à deux, d'où une infinité de classes de repr. lin. de dim 2. En car. 0, il n'y en a que deux, au contraire.

% archive sheet number: 185
